\begin{examplebox}
\textbf{\textcolor{CExample}{例 1（基础）}}

\textbf{\textcolor{CExample}{题目：}} 已知 $x>0, y>0, z>0$，且 $x+y+z=6$，求 $xyz$ 的最大值。

\textbf{\textcolor{CExample}{解题步骤：}}
\begin{description}[style=nextline, leftmargin=2.8em, labelsep=0.5em, itemsep=0.45em]
    \item[\textbf{第一步（应用三元均值）：}] 对正数 $x,y,z$ 直接应用三元均值不等式。
    \[
    \frac{x+y+z}{3} \ge \sqrt[3]{xyz}
    \]
    \par\textbf{理由：} 题目满足“一正”（$x,y,z>0$），且和为定值，符合应用条件。

    \item[\textbf{第二步（代入定值）：}] 将已知条件 $x+y+z=6$ 代入不等式左边。
    \[
    \frac{6}{3} \ge \sqrt[3]{xyz} \quad \Rightarrow \quad 2 \ge \sqrt[3]{xyz}
    \]
    \par\textbf{理由：} 这是简单的算术运算。

    \item[\textbf{第三步（解出最大值）：}] 对不等式 $2 \ge \sqrt[3]{xyz}$ 两边同时立方，得到 $xyz$ 的上界。
    \[
    xyz \le 2^3 = 8
    \]
    \par\textbf{理由：} 立方运算在非负实数上单调递增，不等号方向不变。

    \item[\textbf{第四步（验证取等）：}] 检查等号成立条件。均值不等式等号成立当且仅当 $x=y=z$。结合 $x+y+z=6$，可得 $x=y=z=2$。
    \[
    \text{当 } x=y=z=2 \text{ 时，} xyz = 2 \times 2 \times 2 = 8
    \]
    \par\textbf{理由：} 存在这样一组正数使得等号成立，因此求出的上界 $8$ 是能够取到的最大值。
\end{description}

\textbf{\textcolor{CExample}{关键结论：}} \textbf{在 $x+y+z=6$ 的条件下，$xyz$ 的最大值为 $8$，当且仅当 $x=y=z=2$ 时取得。}

\medskip
\textbf{\textcolor{CExample}{例 2（稍变形）}}

\textbf{\textcolor{CExample}{题目：}} 已知 $a>0, b>0, c>0$，求证：$(a+b+c)\left(\frac{1}{a}+\frac{1}{b}+\frac{1}{c}\right) \ge 9$。

\textbf{\textcolor{CExample}{解题步骤：}}
\begin{description}[style=nextline, leftmargin=2.8em, labelsep=0.5em, itemsep=0.45em]
    \item[\textbf{第一步（分别应用均值）：}] 对 $a,b,c$ 和 $\frac{1}{a},\frac{1}{b},\frac{1}{c}$ 分别应用三元均值不等式。
    \[
    \frac{a+b+c}{3} \ge \sqrt[3]{abc}, \quad \frac{\frac{1}{a}+\frac{1}{b}+\frac{1}{c}}{3} \ge \sqrt[3]{\frac{1}{abc}}
    \]
    \par\textbf{理由：} 因为 $a,b,c>0$，所以 $\frac{1}{a}, \frac{1}{b}, \frac{1}{c}$ 也大于 $0$，满足正数条件。

    \item[\textbf{第二步（改写不等式）：}] 将上一步得到的两个不等式分别变形。
    \[
    a+b+c \ge 3\sqrt[3]{abc}, \quad \frac{1}{a}+\frac{1}{b}+\frac{1}{c} \ge 3\sqrt[3]{\frac{1}{abc}} = \frac{3}{\sqrt[3]{abc}}
    \]
    \par\textbf{理由：} 这是简单的移项和化简，目的是得到题目中左式两部分的下界。

    \item[\textbf{第三步（两式相乘）：}] 将得到的两式左边相乘，右边相乘。
    \[
    (a+b+c)\left(\frac{1}{a}+\frac{1}{b}+\frac{1}{c}\right) \ge \left(3\sqrt[3]{abc}\right) \times \left(\frac{3}{\sqrt[3]{abc}}\right) = 9
    \]
    \par\textbf{理由：} 两个正数不等式（不等号同向）可以相乘，不等号方向不变。

    \item[\textbf{第四步（验证取等）：}] 等号成立的条件是两个均值不等式同时取等，即 $a=b=c$ 且 $\frac{1}{a}=\frac{1}{b}=\frac{1}{c}$，这同时成立当且仅当 $a=b=c$。
    \par\textbf{理由：} 检查取等条件是否可达，以确保证明严密。当 $a=b=c$ 时，左边 $= (3a)(\frac{3}{a})=9$，等号成立。
\end{description}

\textbf{\textcolor{CExample}{关键结论：}} \textbf{已证明对于任意正实数 $a,b,c$，恒有 $(a+b+c)(\frac{1}{a}+\frac{1}{b}+\frac{1}{c}) \ge 9$，且当 $a=b=c$ 时取等号。}
\end{examplebox}
